\clearpage\fancyhead[L]{\small GAUSS AN VON LINDENAU.}\hypertarget{native-331-lower}{}\label{native-331-lower}
\begin{center}\textbf{GAUSS AN VON LINDENAU.}\par\vspace{4pt}\hrule\vspace{4pt}\fontsize{10}{12}\selectfont Monatliche Correspondenz. Band XXIV. Seite 180\,..\,182. 1811. August.\par\vspace{4pt}\hrule\vspace{4pt}\end{center}

\begin{flushright}Göttingen, den 3. August 1811.\end{flushright}
\Needspace{16\baselineskip}\par\indent Ich habe in den letzten Tagen aus den Beobachtungen des Kometen, die Herr \textsc{von Zach} angestellt hat und die Sie an Professor \textsc{Harding} geschickt hatten, vorläufig die folgenden parabolischen Elemente abgeleitet.
\SourceLeaf{332}
\begin{center}\fontsize{10}{12}\selectfont Durchgang durch das Perihel 1811 Sept. 10\\$0^{\text{u}}\;51^{\text{m}}\;21^{\text{s}}$ Merid. von Göttingen.\end{center}

\begin{center}\fontsize{10}{12}\selectfont\setlength{\tabcolsep}{4pt}\renewcommand{\arraystretch}{1.15}\begin{tabular}{llr}
Logarithmus des Abstands im Perihel & .\;.\;.\;.\;.\;.\;. & $9.99153$\\
Länge des Perihels & .\;.\;.\;.\;.\;.\;.\;.\;.\;.\;.\;.\;.\;.\;. & $73^\circ\;14'\;35''$\\
Länge des aufsteigenden Knotens & .\;.\;.\;.\;.\;.\;.\;.\;.\;.\;.\;. & $141\;4\;59$\\
Neigung der Bahn & .\;.\;.\;.\;.\;.\;.\;.\;.\;.\;.\;.\;.\;.\;.\;.\;. & $73\;48\;2$\\
Bewegung rückläufig &  & \\
\end{tabular}\end{center}

\Needspace{15\baselineskip}\par\indent Nach der Mitte dieses Monats kann man bereits damit beginnen, diesen bemerkenswerten Himmelskörper am nördlichen Himmel zu suchen. Dafür habe ich aus den obigen Elementen die folgenden Positionen berechnet.

\begin{center}\fontsize{10}{12}\selectfont\setlength{\tabcolsep}{4pt}\renewcommand{\arraystretch}{1.15}\begin{tabular}{lrr}
1811 Aug. 11 $12^{\text{u}}$ & AR. $142^\circ\;43'$ & Decl. $29^\circ\;54'$ N.\\
26 & $152\;19$ & $36\;15$\\
Sept. 10 & $167\;34$ & $43\;14$\\
\end{tabular}\end{center}

\begin{center}\fontsize{10}{12}\selectfont\setlength{\tabcolsep}{4pt}\renewcommand{\arraystretch}{1.15}\begin{tabular}{lrr}
1811 Sept. 25 $12^{\text{u}}$ & AR. $194^\circ\;22'$ & Decl. $48^\circ\;35'$ N.\\
Oct. 10 & $232\;28$ & $44\;52$\\
25 & $262\;49$ & $30\;42$\\
\end{tabular}\end{center}

\Needspace{20\baselineskip}\par\indent Diese Ergebnisse können jedoch nur als vorläufig gelten, weil die bisher langsame heliozentrische Bewegung (vom 11. April bis zum 2. Juni nur $15^\circ$) noch keine sehr genaue Bestimmung der Elemente erlaubt. Die Helligkeit, bezogen auf die Helligkeit im Abstand 1 von Erde und Sonne als Einheit, hatte folgende Werte

\begin{center}\fontsize{10}{12}\selectfont\setlength{\tabcolsep}{4pt}\renewcommand{\arraystretch}{1.15}\begin{tabular}{llr}
April 11 & .\;.\;.\;.\;. & $0.035$\\
Mai 7 & .\;.\;.\;.\;. & $0.042$\\
Juni 2 & .\;.\;.\;.\;. & $0.052$\\
\end{tabular}\end{center}

\noindent und wird folgende Werte haben

\begin{center}\fontsize{10}{12}\selectfont\setlength{\tabcolsep}{4pt}\renewcommand{\arraystretch}{1.15}\begin{tabular}{llr}
Aug. 11 & .\;.\;.\;.\;.\;.\;.\;.\;. & $0.194$\\
26 & .\;.\;.\;.\;.\;.\;.\;.\;. & $0.294$\\
Sept. 10 & .\;.\;.\;.\;.\;.\;.\;.\;. & $0.434$\\
Sept. 25 & .\;.\;.\;.\;.\;.\;.\;.\;. & $0.576$\\
Oct. 10 & .\;.\;.\;.\;.\;.\;.\;.\;. & $0.595$\\
25 & .\;.\;.\;.\;.\;.\;.\;.\;. & $0.415$\\
\end{tabular}\end{center}

\noindent Erst gegen Februar wird der Komet im Wassermann unsichtbar. Vielleicht wird er dort im April oder Mai wieder sichtbar, falls sein Licht nicht zu schwach ist. Für den 2. Juli 1812 habe ich die folgenden Werte überschlägig berechnet.
\begin{center}AR. $336^\circ$\quad Dekl. $22^\circ$ südl.\quad Helligkeit $0.005$.\end{center}
\par\indent [Handschriftliche Bemerkung. Nach verbesserten Elementen 1812 Juni $3.22^{\mathrm u}$ $337^\circ\,49'$ --- $11^\circ\,41'$ --- $0.5435$ $0.0057$.]

\par\indent Ich habe die Berechnungen der Störungen von \emph{Pallas} nun bis Ende 1812 fortgesetzt. In einigen Tagen hoffe ich Ihnen die Ephemeride schicken zu können, die Herr \textsc{Nicolai}, ein sehr geübter junger Mann, gerade berechnet.
